% ============================================================================
% Geopolitical Oil Shocks, Sectoral Heterogeneity, and Aggregation Masking:
% Evidence from a Multi-Event Framework
% Target Journal: Research in International Business and Finance (RIBAF)
% ============================================================================

\documentclass[12pt,preprint,authoryear]{elsarticle}

% ---- Packages ----
\usepackage[utf8]{inputenc}
\usepackage[T1]{fontenc}
\usepackage{amsmath,amssymb}
\usepackage{graphicx}
\usepackage{booktabs}
\usepackage{multirow}
\usepackage{setspace}
\usepackage{hyperref}
\usepackage[margin=1in]{geometry}
\usepackage{natbib}
\usepackage{caption}
\usepackage{subcaption}
\usepackage{threeparttable}
\usepackage{float}
\usepackage{dcolumn}
\usepackage{tabularx}
\usepackage{pdflscape}
\usepackage{array}

% ---- Settings ----
\doublespacing
\hypersetup{colorlinks=true,citecolor=blue,linkcolor=blue,urlcolor=blue}

% ---- Journal metadata ----
\journal{Research in International Business and Finance}

\begin{document}

\begin{frontmatter}

\title{Geopolitical Oil Shocks, Sectoral Heterogeneity, and Aggregation Masking: Evidence from a Multi-Event Framework}

\author[inst1]{Apoorv Saxena}
\ead{s1129420@mail.yzu.edu.tw}

\author[inst1]{Yan-Jie Yang\corref{cor1}}
\ead{yanjie@saturn.yzu.edu.tw}

\cortext[cor1]{Corresponding author.}

\affiliation[inst1]{organization={College of Management, Yuan Ze University},
                     city={Taoyuan},
                     country={Taiwan}}

\begin{abstract}
Financial economists and practitioners routinely rely on aggregate market indices as summary statistics of economic conditions and the impact of macroeconomic shocks. This paper challenges the assumption that such indices accurately reflect the underlying economic effects of geopolitical oil shocks. Using a multi-event framework spanning five major oil-related shocks---the 2026 Strait of Hormuz crisis, the Russia--Ukraine conflict, OPEC+ production cuts, the Hamas attack, and the COVID-19 demand collapse---we show that pre-event oil price sensitivity ($\beta_{\text{oil}}$) strongly predicts cross-sectional cumulative abnormal returns ($\gamma_1 = +20.18$, $p < 0.0001$ with GICS-clustered standard errors). The transmission channel operates asymmetrically: positively for supply-side shocks ($\gamma_1 = +19.50$) and negatively for the demand-side shock ($\gamma_1 = -32.61$), with an interaction coefficient of $+58.79$ ($p < 0.0001$). To quantify information loss in aggregate indices, we introduce a bounded Aggregation Masking Index ($\Psi$), which reveals that approximately 94.7\% of sectoral return variation is concealed by capitalization-weighted aggregation across all five events. These findings demonstrate that reliance on aggregate indices systematically understates the economic impact of geopolitical shocks, with implications for policymakers, investors, and researchers who use market indices to infer macroeconomic conditions.
\end{abstract}

\begin{keyword}
Geopolitical risk \sep Oil price shocks \sep Sectoral stock returns \sep Aggregation bias \sep Event study \sep Multi-event panel
\end{keyword}

\begin{keyword}[JEL]
G14 \sep G12 \sep F51 \sep Q43
\end{keyword}

\end{frontmatter}

% ====================================================================
% 1. INTRODUCTION
% ====================================================================
\section{Introduction}

Financial economists and practitioners routinely rely on aggregate market indices---such as the S\&P~500---as summary statistics of economic conditions and the impact of macroeconomic shocks. This widespread practice implicitly assumes that aggregate indices accurately reflect the underlying economic effects of such shocks. In this paper, we challenge this assumption. We argue that when shocks generate heterogeneous responses across sectors, aggregate indices can systematically misrepresent the true economic impact.

Geopolitical shocks provide a particularly relevant setting in which this issue arises. Such shocks are typically transmitted through specific economic channels---most notably energy prices, supply chains, and trade exposure---affecting industries asymmetrically \citep{kilian2009not,caldara2022measuring}. As a result, some sectors benefit while others experience substantial losses. When these heterogeneous responses are aggregated through capitalization-weighted indices, positive and negative effects may offset each other, leading to a muted aggregate response. Consequently, aggregate indices may exhibit a form of aggregation bias, obscuring economically meaningful sectoral disruptions.

We study this phenomenon using a multi-event framework centered on five major geopolitical oil shocks: the 2026 Strait of Hormuz crisis, the 2022 Russia--Ukraine conflict, the 2022 OPEC+ production cut, the 2023 Hamas attack, and the 2020 COVID-19 demand collapse. These events span both supply-side and demand-side oil shocks, enabling us to test whether the transmission channel operates differently across shock types---a distinction emphasized by \citet{kilian2009not} and \citet{kilian2009impact}.

We exploit the February 28, 2026 Strait of Hormuz disruption as our primary quasi-natural experiment. The Strait handles approximately 21\% of global petroleum consumption daily \citep{eia2023hormuz}, and the crisis was sudden, largely unanticipated, and primarily transmitted through oil price channels, making it ideal for cross-sectional identification. Using cross-sectional regressions on 34 sector and sub-sector ETFs, we show that pre-event oil price sensitivity ($\beta_{\text{oil}}$) strongly predicts cumulative abnormal returns (CARs), with $\gamma_1 = +20.18$ ($p < 0.0001$) that survives six increasingly saturated specifications including market risk, macroeconomic factors, and firm fundamentals. This finding generalizes across all five events in our panel.

To quantify the extent to which aggregate indices mask these heterogeneous responses, we introduce a novel bounded measure---the Aggregation Masking Index ($\Psi$). Unlike simple ratios of index-level to sectoral returns, $\Psi$ is bounded on $[0,1]$ by Jensen's inequality and captures the degree to which capitalization-weighted aggregation conceals underlying sectoral variation. Empirically, we find $\Psi \approx 0.947$ across all five events, implying that approximately 94.7\% of economically meaningful sectoral variation is hidden through aggregation. Crucially, this result holds regardless of shock type or magnitude, suggesting that aggregation masking is a structural feature of capitalization-weighted indices rather than an event-specific phenomenon.

This paper makes four contributions. First, we demonstrate that aggregate market indices systematically underestimate the economic impact of geopolitical shocks when sectoral responses are heterogeneous, contributing to the literature on aggregation and measurement in asset pricing \citep{vuolteenaho2002what}. Second, we identify oil price exposure as a key transmission mechanism linking geopolitical shocks to cross-sectional equity returns, extending the oil shock heterogeneity literature \citep{kilian2009not,arouri2012oil,broadstock2012oil}. Third, we propose $\Psi$ as a tractable and interpretable metric for quantifying aggregation-induced information loss, providing a new tool for measuring the gap between index-level and sector-level responses. Fourth, we establish that the oil-beta transmission channel operates asymmetrically across shock types---positively for supply-side shocks and negatively for demand-side shocks---elevating the analysis from a single-event case study to a generalizable framework.

The remainder of the paper is organized as follows. Section~\ref{sec:background} provides institutional background and develops the hypotheses. Section~\ref{sec:data} describes the data and methodology. Section~\ref{sec:results} presents the main results on heterogeneous shock transmission and mechanism identification. Section~\ref{sec:external} extends the analysis to a multi-event panel and tests shock-type asymmetry. Section~\ref{sec:masking} analyzes aggregation masking. Section~\ref{sec:conclusion} concludes.

% ====================================================================
% 2. EVENT BACKGROUND AND HYPOTHESES DEVELOPMENT
% ====================================================================
\section{Event background and hypotheses development}\label{sec:background}

\subsection{The 2026 Strait of Hormuz crisis}

On February 28, 2026, escalating tensions between Iran and a US-Israeli coalition culminated in military strikes near the Strait of Hormuz, the world's most critical oil transit chokepoint. The Strait connects the Persian Gulf to the Gulf of Oman and handles approximately 21 million barrels per day---roughly one-fifth of global petroleum consumption \citep{eia2023hormuz}. Iran's subsequent announcement of a partial blockade on March 3, followed by the first physical interception of a commercial tanker on March 5, triggered an immediate spike in crude oil prices and widespread uncertainty across global financial markets.

This event provides a particularly clean identification setting for several reasons. First, the shock was sudden and largely unanticipated in its severity, satisfying the exogeneity condition required for event study methodology \citep{mackinlay1997event}. Second, the primary transmission channel---oil prices---is clearly identified and directly measurable, enabling us to construct sector-specific exposure measures. Third, the event generated sharp sectoral divergence: oil-linked ETFs such as USO gained over 42\%, while oil-cost-sensitive sectors such as homebuilders (ITB) lost over 16\%, even as the S\&P~500 declined by only 1.81\%.

\subsection{Additional events and shock classification}

To establish external validity and test shock-type asymmetry, we extend the analysis to four additional major oil-related geopolitical events:

\begin{itemize}
    \item \textbf{Russia--Ukraine conflict} (February 24, 2022): A supply-side shock that disrupted Russian oil and gas exports to Europe, triggering energy price surges and sanctions-related supply concerns.
    \item \textbf{OPEC+ production cut} (October 5, 2022): A coordinated supply restriction of 2 million barrels per day, representing a deliberate supply-side intervention.
    \item \textbf{Hamas attack} (October 9, 2023): A geopolitical shock raising concerns about broader Middle Eastern supply disruptions, classified as a supply-side risk event.
    \item \textbf{COVID-19 demand collapse} (March 9, 2020): A demand-side shock driven by the global pandemic that caused an unprecedented collapse in oil consumption and prices.
\end{itemize}

Following \citet{kilian2009not}, we classify the first four events as supply-side shocks and COVID-19 as a demand-side shock. This classification is central to our fourth hypothesis: supply shocks create sectoral winners and losers (positive $\gamma_1$), while demand shocks cause correlated losses across all sectors, with high-oil-beta sectors falling disproportionately (negative $\gamma_1$).

\subsection{Hypotheses development}

Geopolitical shocks are often transmitted through specific economic channels, with oil prices playing a central role \citep{kilian2009not,hamilton2003oil}. However, sectors differ in their exposure to oil price fluctuations, leading to heterogeneous responses in equity returns \citep{arouri2012oil,scholtens2012oil}. When these heterogeneous responses are aggregated into market indices, offsetting effects may obscure the true economic impact of shocks. Based on this framework, we develop the following hypotheses.

\bigskip
\noindent\textbf{H1: Heterogeneous Shock Transmission.} \textit{The impact of geopolitical oil shocks on equity returns is heterogeneous across sectors and is systematically related to sectoral exposure to oil price fluctuations.}

We test this hypothesis using cross-sectional regressions of the form:
\begin{equation}\label{eq:h1}
    \text{CAR}_i = \alpha + \gamma_1 \beta_{\text{oil},i} + \boldsymbol{\gamma}' \mathbf{X}_i + u_i
\end{equation}
where $\beta_{\text{oil},i}$ is the pre-event oil price sensitivity and $\mathbf{X}_i$ is a vector of controls. H1 predicts $\gamma_1 \neq 0$.

\bigskip
\noindent\textbf{H2: Mechanism Hypothesis.} \textit{The effect of oil price exposure on sectoral returns is amplified in sectors with higher energy intensity, supply-chain dependence, and trade exposure.}

We test this using interaction regressions:
\begin{equation}\label{eq:h2}
    \text{CAR}_i = \alpha + \gamma_1 \beta_{\text{oil},i} + \gamma_2 M_i + \gamma_3 (\beta_{\text{oil},i} \times M_i) + \boldsymbol{\gamma}' \mathbf{X}_i + u_i
\end{equation}
where $M_i$ represents mechanism variables (energy intensity, supply chain dependence, or trade exposure). H2 predicts $\gamma_3 > 0$ for energy intensity and supply chain, indicating that these channels amplify the oil transmission mechanism.

\bigskip
\noindent\textbf{H3: Aggregation Masking Hypothesis.} \textit{When sectoral responses to geopolitical shocks are heterogeneous, aggregate market indices systematically understate the magnitude of economic disruption.}

We test this using the Aggregation Masking Index:
\begin{equation}\label{eq:h3}
    \Psi_e = 1 - \frac{\left|\sum_{i=1}^{N} w_i \cdot \text{CAR}_{i,e}\right|}{\sum_{i=1}^{N} w_i \cdot |\text{CAR}_{i,e}|}
\end{equation}
where $w_i$ are capitalization weights. By Jensen's inequality, $\Psi \in [0,1]$. H3 predicts $\Psi > 0$ and substantively large, indicating meaningful information loss through aggregation.

\bigskip
\noindent\textbf{H4: Shock-Type Asymmetry.} \textit{The oil-beta transmission channel operates asymmetrically across shock types: positively for supply-side shocks and negatively for demand-side shocks.}

We test this in a multi-event panel:
\begin{equation}\label{eq:h4}
    \text{CAR}_{i,e} = \beta_1 \beta_{\text{oil},i} + \beta_2 (\beta_{\text{oil},i} \times \text{Supply}_e) + \mu_e + \varepsilon_{i,e}
\end{equation}
where $\text{Supply}_e$ is a dummy equal to one for supply-side events and $\mu_e$ are event fixed effects. H4 predicts $\beta_2 > 0$, indicating that the total oil-beta effect ($\beta_1 + \beta_2$) is positive for supply shocks but negative ($\beta_1$ alone) for demand shocks.

% ====================================================================
% 3. DATA AND METHODOLOGY
% ====================================================================
\section{Data and methodology}\label{sec:data}

\subsection{Data}

We collect daily returns for 37 US sector and sub-sector ETFs from January 2018 to March 25, 2026 via Yahoo Finance. We apply a liquidity filter requiring average daily dollar volume exceeding \$10 million. After dropping one ETF for insufficient liquidity (HACK) and excluding SPY (used as the market benchmark), the cross-sectional sample comprises 34 ETFs spanning all 11 GICS sectors and 23 sub-industry ETFs.\footnote{The ETF universe includes GICS sector ETFs (XLK, XLF, XLV, XLE, XLI, XLY, XLP, XLRE, XLU, XLB, XLC) and sub-sector ETFs (USO, XOP, FCG, OIH, AMLP, TAN, URA, SMH, IGV, IBB, XBI, IHF, KBE, KRE, KIE, IAI, ITA, XAR, ITB, XHB, IYT, JETS, XRT, REMX). SPY serves as the market return proxy.}

Oil prices (WTI crude), the CBOE Volatility Index (VIX), the US Dollar Index (DXY), and 10-year Treasury yields are obtained from FRED. Sector-level mechanism variables---energy intensity, supply chain dependence, and trade exposure---are constructed from the Bureau of Economic Analysis (BEA) 2022 Annual Input-Output Tables (Use Table and Total Requirements Table) and Census Bureau trade data \citep{bea2022io}.

\subsection{Event study methodology}

We estimate market-model abnormal returns following \citet{mackinlay1997event}:
\begin{equation}\label{eq:mktmodel}
    R_{it} = \alpha_i + \beta_i R_{mt} + \varepsilon_{it}
\end{equation}
where the estimation window spans 240 trading days ending 10 days before the event. For the primary Hormuz event, this corresponds to March 3, 2025 to February 12, 2026. Cumulative abnormal returns are computed over the primary $[-10, +10]$ window:
\begin{equation}
    \text{CAR}_i = \sum_{t=-10}^{+10} \hat{\varepsilon}_{it}
\end{equation}
Statistical significance is assessed using \citet{kolari2010event} adjusted $t$-statistics, which account for cross-sectional correlation and event-induced volatility.

\subsection{Multi-factor beta estimation}

To isolate oil price sensitivity from correlated factor exposures, we estimate a five-factor model in the pre-event estimation window:
\begin{equation}
    R_{it} = \alpha_i + \beta_{\text{mkt},i} R_{mt} + \beta_{\text{oil},i} R_{\text{oil},t} + \beta_{\text{vix},i} \Delta\text{VIX}_t + \beta_{\text{dxy},i} R_{\text{dxy},t} + \beta_{\text{y10},i} \Delta Y_{10,t} + \varepsilon_{it}
\end{equation}
where $R_{\text{oil},t}$ is the daily WTI crude oil return. The resulting $\beta_{\text{oil},i}$ measures each ETF's sensitivity to oil price movements, conditional on market, volatility, currency, and interest rate factors.

\subsection{Cross-sectional regression specifications}

We estimate six nested specifications of increasing saturation:
\begin{equation}\label{eq:cross}
    \text{CAR}_i = \alpha + \gamma_1 \beta_{\text{oil},i} + \boldsymbol{\gamma}' \mathbf{X}_i + u_i
\end{equation}

Specification~(1) includes only $\beta_{\text{oil}}$. Specification~(2) adds $\beta_{\text{mkt}}$. Specifications~(3)--(4) add macroeconomic factor betas ($\beta_{\text{vix}}$, $\beta_{\text{dxy}}$, $\beta_{\text{y10}}$). Specifications~(5)--(6) add fundamental characteristics: historical volatility, log total assets, debt-to-assets ratio, and price-to-book ratio. All standard errors are HC3-robust \citep{long2000using}, with GICS-sector clustering as a robustness check.

\subsection{Mechanism variables}

To test H2, we construct three sector-level mechanism variables from BEA Input-Output data:

\begin{itemize}
    \item \textbf{Energy intensity}: The share of petroleum, coal, and utility inputs in total intermediate consumption, computed from the BEA 2022 Use Table. Values range from 2.1\% (Communication Services) to 82.3\% (Energy).
    \item \textbf{Supply chain dependence}: The direct and indirect petroleum requirements coefficient from the BEA 2022 Total Requirements Table. Values range from 0.03 (Communication Services) to 0.92 (Energy).
    \item \textbf{Trade exposure}: The ratio of exports plus imports to total shipments, computed from Census Bureau USA Trade Online data. Values range from 2.9\% (Utilities) to 52.8\% (Technology).
\end{itemize}

Sub-sector ETFs inherit the mechanism variable values of their parent GICS sector. All mechanism variables are standardized (mean zero, unit variance) before entering interaction terms.

\subsection{Aggregation Masking Index ($\Psi$)}

We define the bounded Aggregation Masking Index as:
\begin{equation}\label{eq:psi}
    \Psi_e = 1 - \frac{\left|\sum_{i=1}^{N} w_i \cdot \text{CAR}_{i,e}\right|}{\sum_{i=1}^{N} w_i \cdot |\text{CAR}_{i,e}|}
\end{equation}
where $w_i$ are capitalization weights that sum to one. By Jensen's inequality, $|\sum w_i \cdot \text{CAR}_i| \leq \sum w_i \cdot |\text{CAR}_i|$, ensuring $\Psi \in [0,1]$. A value of $\Psi = 0$ indicates no masking (all sectors move in the same direction), while $\Psi = 1$ indicates complete masking (the weighted sum of CARs is exactly zero despite nonzero individual CARs). We interpret $\Psi$ as the fraction of sectoral return variation that is concealed through capitalization-weighted aggregation.

We also report the traditional ratio-based index $\Phi = |\text{CAR}_{\text{SPY}}| / \overline{|\text{CAR}_i|}$ for comparison, where values below one indicate that the aggregate index understates average sectoral disruption.

% ====================================================================
% SECTIONS 4-7 CONTINUED IN SUBSEQUENT CHUNKS
% ====================================================================

\end{document}
